Agradeço, antes de tudo, a Deus, por ter me dado saúde e persistência para concluir esta etapa. Agradeço também à minha mãe, Quitéria Alves, à minha querida avó, Cícera Monteiro, e à minha irmã gêmea, Isabel Alessandra. Elas são o meu maior apoio, e tudo o que sou hoje, tudo o que conquistei, devo a elas. Minha avó esteve presente em cada madrugada de estudo. Minha mãe sempre apoiou minhas decisões, incentivando o que me fazia feliz e realizada. Tenho orgulho, também, de dividir com ela a paixão pelo esporte — foi ela quem, desde cedo, me ensinou a torcer e que já foi atleta de futebol feminino; hoje, assistimos juntas aos jogos da NBA, torcendo lado a lado. Não por acaso, essa paixão compartilhada está no coração deste trabalho. E à minha irmã, a minha pessoa favorita no mundo, obrigada por caminhar comigo desde sempre.

Ao meu namorado, Luann Ferreira, meu companheiro e parceiro de todas as horas, agradeço por tudo o que vivemos desde que nos conhecemos no curso. Fomos dupla de laboratório e de tantos outros projetos, e hoje dividimos a vida. Obrigada por estar ao meu lado também nesta etapa.

Ao meu orientador, Professor Dr. Dimas Cassimiro do Nascimento Filho, agradeço pela orientação cuidadosa, pela disponibilidade em discutir cada etapa da pesquisa e pela confiança depositada.

Estendo esse agradecimento aos professores do curso de Ciência da Computação, pelos ensinamentos que tornaram esta pesquisa possível.

Aos amigos que fiz durante a graduação — Ana Beatriz, Maria Virgínia, Thiago Fabrício, Gabriel Viana, Geovânio José Silva, Pedro Vinícius, Victor Hugo, Rodrigo Leandro e Maria Luiza —, obrigada pelas trocas de ideias, pelo companheirismo nas disciplinas mais difíceis e por tornarem esses anos mais leves.

Por fim, a todos que, de alguma forma, contribuíram para a realização deste trabalho.
